

\subsection{Key observations}

Given a \emph{restricted} \laVASS $V$ in dimension $d$, assuming that
$L(V)\neq \emptyset$, there is a run $\pi$ from an initial
configuration $c$ to a final configuration $c'$. With no loss of
generality, throughout this section, we assume that $\val(c',x_i) \ge
\val(c',x_{i+1})>0$ for all $1\le i<d$. In particular, this implies
that every counter gets incremented at least once along a path
witnessing non-emptiness.

Our first observation is that if along $\pi$ a counter $y_i$ achieves
the first time a non-zero value by taking a $\wTTPO$ labeled
transition, the length of the remaining segment of $\pi$ is bounded by
$O(\log(m_i+1))$, where $m_i$ is the value of counter $x_i$ before the
transition is taken. The reason is that, once $y_i$ has non-zero
value, its value at least doubles when a transition is taken. Hence if
$\pi$ is ``long'' then along $\pi$ there will be loops incrementing a
counter $x_i$ before the corresponding $y_i$ achieves non-zero value.

In the latter scenario, we may actually, subject to some bookkeeping,
discard concrete values of $x_i$ and $y_i$ and only store their
residue classes modulo $\ell_i$, where $\ell_i$ is the length of the
first loop incrementing $x_i$ along $\pi$. In particular, if we are
given a non-accepting run $\pi'$ such that $\val(\pi',x_i) \equiv
\val(\pi',y_i) \bmod \ell_i$ and $\val(\pi',x_i)<\val(\pi',y_i)$ then
$\pi'$ can be turned into a run $\pi''$ where $\val(\pi'',x_i) =
\val(\pi'',y_i)$ by iterating the loop of length $\ell_i$.

There are, however, some further subtleties that need to be taken care
of. Consider the segment $\pi'$ of $\pi$ between the first transition
labeled by $\wPP$ on $x_i$ and the first transition labeled by $\wPP$ on
$x_{i+1}$. If $\pi'$ contains no loop then we are in a situation where
the first loop incrementing $x_i$ is also the first loop incrementing
$x_{i+1}$. This means that the values of $x_i$ and $x_{i+1}$ get
paired together, and hence, for an accepting run, also the values of
$y_i$ and $y_{i+1}$ are paired together. In our approach, we deal with
such circumstances by introducing so-called \emph{$y$-constraints}. A
$y$-constraint of the form $y_i-y_{i+1}=\delta_i$ for some constant
$\delta_i \in \N$ asserts that the counters $y_{i+1}$ and $y_i$ must
eventually have constant difference $\delta_i$ along a run.

Otherwise, if $\pi'$ above contains a loop, the difference between the
values of $x_i$ and $x_{i+1}$ is not necessarily constant, but
lower-bounded by the length $\delta_i$ of the loop-free sub path of
$\pi'$. Thus, in an accepting run, the difference between $y_i$ and
$y_{i+1}$ must also be at least $\delta_i$, which is asserted by a
$y$-constraint of the form $y_i-y_{i+1} \ge \delta_i$.



%% \subsection{Shortest accepting paths can be very long}

%% We present an example of a restricted \laVASS that has an accepting
%% path, but the length of the shortest such path is non-elementary. All
%% the unspecified transitions labels are $\timestwo$ for the y-counters
%% and $id$ for the x-counters with value zero, and $\wPP$ for the counters
%% with value greater than zero, respectively.

%% \begin{figure}[!t]

%% \begin{center}
%% \begin{tikzpicture}[->,shorten >=1pt,auto,node distance=2.2cm,semithick]

%% \tikzstyle{every state}=[draw=black,text=black, minimum size = 1mm]

%% \node[initial by arrow,state,initial text=] (1) {};
%% \node[state] (2) [right of=1] {};

%% \node[state] (4) at (0,-1.1) {};
%% \node[state] (3) [right of=4] {}; 
%% \node[state] (5) [right of=3] {};
%% \node[state, accepting] (9) [right of=5] {};

%% \path 
%%     (1) edge node [align=center]{$y_d:\timestwob$}  (2)
%%     (1) edge [loop above] node {$x_d:\wPP$}(1)
    
%%     (2) edge [loop above] node {$x_{d-1}:\wPP$} (2)
    
%%     (2) edge [dotted] node {} (4)


%%     (4) edge node [align=center,below]{$y_{2}:\timestwob$}  (3)
%%     (4) edge [in=120,out=60,loop below] node {$x_{2}:\wPP$}(4)
    
%%     (3) edge node [align=center,below]{$y_1:\timestwob$}  (5)
%%     (3) edge [in=120,out=60,loop below] node {$x_1:\wPP$}(3)
    
    
%%     (5) edge node [align=center,below]{$y_1:\timestwob$}  (9);
    

    
%% \end{tikzpicture}
%% \end{center}

%% \caption{An \laVASS with an accepting path requiring non-elementary
%%   length.}
%% \label{fig:non-elem}
%% \end{figure}

%% Let $c=(q,m_1,n_1,\dots,m_d,n_d)$ be the final configuration we reach after we take the shortest accepting path of the \laVASS in \Cref{fig:non-elem}. Clearly, $n_1$ (the value of the counter $y_1$) is 3 and so, $m_1$ must be 3 since $c$ is a final configuration. 

%% Now assume that $m_i,n_i\geq2\uparrow\uparrow i$ (here, $\uparrow$ is the Knuth up-arrow notation), for all $i\in\{1,\dots,k-1\}$. Let us consider the values of $m_{k},n_k$:

%% \begin{figure}[h]

%% \begin{tikzpicture}[->,shorten >=1pt,auto,node distance=2.2cm,semithick]

%% \tikzstyle{every state}=[draw=black,text=black, minimum size = 3mm]

%% \node (1) at (2.5,0) {.};
%% \node (2) at (3,0) {.};
%% \node (3) at (3.5,0) {.};
%% \node[state] (4) at (4.5,0) {};
%% \node[state] (5) [right of=4] {}; 
%% \node (6) at (8,0) {.};
%% \node (7) at (8.5,0) {.};
%% \node (8) at (9,0) {.};

%% \path 

%%     (4) edge node [align=center]{$y_{k}:\timestwob$}  (5)
%%     (4) edge [in=120,out=60,loop above] node {$x_{k}:\wPP$}(4)
    
%%     (5) edge [in=120,out=60,loop above] node {$x_{k-1}:\wPP$}(5);
    
    
%% \end{tikzpicture}
%% \centering

%% \label{fig:induction step}
%% \end{figure}

%% Since we know that $m_{k-1}\geq 2\uparrow\uparrow (k-1)$ it means that $n_k\geq 2^{2\uparrow\uparrow (k-1)}=2\uparrow\uparrow k$. However, we know that $m_k=n_k$, so $m_k\geq 2\uparrow\uparrow k$. The key observation is that $x_{k-1}$ needs to take many self-loops to recover the difference to the $y_{k-1}$ counter. However, every time $x_{k-1}$ takes a self-loop, the value of $y_{k}$ doubles and $x_k$ needs to take an exponential number of self-loops in order to recover that, which in turn doubles the value of counter $y_{k+1}$ each time. 

%% We can conclude that $m_d\geq 2\uparrow\uparrow d$ and since $m_d$ is
%% equal to the length of the shortest accepting path (we cannot take any
%% transition in the beginning without updating the value of $m_1$), we
%% have shown that the length of the shortest accepting path of a \laVASS
%% can be non-elementary with respect to the size of the input restricted
%% \laVASS.



\subsection{An abstraction for restricted \laVASS}

Our decision procedure for emptiness of restricted \laVASS is based on
reducing this problem to a reachability problem in a
carefully designed finite-state abstraction of the state-space of
\laVASS. Throughout this section, let $V=\langle Q,2d,
\Sigma,\Delta,\lambda,q_0,F,\Phi\rangle$ be a restricted \laVASS.  We
first define the state space of the abstracted \laVASS.
%
\begin{definition}
  An abstract configuration is a tuple
%
  \begin{multline*}
  \alpha=(q,
  m_1,n_1,\dots,m_{d},n_{d},u_1,u_2,\dots,u_{d-1},\ell_1,\dots,\ell_d)\\
  \in Q\times (\N\cup\{\bot\})^{2d} \times \N^{d-1}\times(\N \cup \{ \top \})^{d},
  \end{multline*}
  %
  such that $m_i,n_i \in [0,2dM_i] \cup \{ \bot \}$
  and $u_i \in [0,U_i]$ and $\ell_i\in[0,M_i-1] \cup \{\top\}$,
  where
  \begin{itemize}
	\item $M_i \defeq \floor{\abs Q^{((1/8)\cdot 32^{i-1} + 1)}}$; and
	\item $U_i \defeq \abs Q^{(32^{i-1}+4)}$.
  \end{itemize}
\end{definition}
%
The idea is that $m_i,n_i$ store the residue classes modulo $\ell_i$
of the counter pair $x_i,y_i$ respectively, where the value $\top$ for
$\ell_i$ acts as an indicator that we are storing actual values and
not residue classes. The value $\bot$ for some $x_i$ or $y_i$
indicates that the counter has not yet been initialised. If for an
update function $f$, $f=\wPP$ or $f=\wTTPO$ then $f(\bot)\defeq
1$; otherwise $f(\bot) \defeq \bot$, and we stipulate that $\bot \bmod
n = \bot$. The value of $u_i$ in an abstract configuration carries the
current difference between the value of the counters $y_{i}$ and
$y_{i+1}$. This difference is potentially unbounded, however for our
purposes it suffices to only store its value if it is less than $U_i$,
and to indicate the fact that it is at least $U_i$ by the value
$u_i=U_i$.
%This difference must not become negative since $u_1,\dots,u_{d-1}$ are
%all natural numbers.

We denote the (finite) set of all abstract configurations of $V$ by
$A(V)$. Let us now define a transition relation
${\xrightarrow{\cdot}}\subseteq A(V) \times \Delta \times A(V)$ such
that $\alpha\xrightarrow{t}\alpha'$, $t=(q,a,q')\in \Delta$ if and
only if:
%
\begin{itemize}
    \item $\alpha=(q,m_1,n_1,\dots,u_{d-1},\ell_1,\dots,\ell_d)$ and $\alpha'=(q',m_1',n_1',\dots,u_{d-1}',\ell_1,\dots,\ell_d)$;
    \item $\lambda(t)=(f_{x_1},f_{y_1},\dots,f_{x_d},f_{y_d})$;
    \item $f_{x_i} = \wPP$ for all $i$ such that $m_i\neq \bot$;
    \item if $\ell_i\neq\top$, $m_i' = f_{x_i}(m_i)\bmod \ell_i$ and $n_i'=f_{y_i}(n_i)\bmod \ell_i$;
    \item if $\ell_i=\top$, $m_i' = f_{x_i}(m_i)$ and $n_i'=f_{y_i}(n_i)$; and
    \item for all $i\in\{1,\dots,d-1\}$,
      $$u_i' =
    \begin{cases*}
      min(2u_i+1,U_i) & if $f_{y_{i}}=\wTTPO, f_{y_{i+1}}=\wTT$ \\
      min(2u_i,U_i) & if $f_{y_{i}}=f_{y_{i+1}}$ \\
      min(2u_i-1,U_i) & if $f_{y_{i}}=\wTT, f_{y_{i+1}}=\wTTPO$. \\
    \end{cases*}
    $$
\end{itemize}
%
Assuming that the value of $y_i$ is at least the value of $y_{i+1}$,
which we will always ensure, the definition of how to update $u_i$
ensures that it exactly stores the difference $y_i-y_{i+1}$ unless the
difference becomes too large, in which case it is levelled off at
$U_i$.

An abstract configuration path is a sequence of abstract
configurations and transitions of the form
$R=\alpha_1\xrightarrow{t_1} \alpha_2 \xrightarrow{t_2} \cdots
\xrightarrow{t_{n-1}} \alpha_n$.

Given two consecutive $y$-counters $y_i,y_{i+1}$ and $\delta_i\in \N$,
we say that $y_{i}-y_{i+1}=\delta_i$ and $y_{i}-y_{i+1}\geq \delta_i$
are \emph{$y$-constraints}. Let $Y$ be a set of $y$-constraints, an
abstract configuration
$\alpha=(q,m_1,n_1,\dots,u_1,\dots,u_{d-1},\ell_1,\dots,\ell_d)$
\emph{respects $Y$} whenever
%
\begin{itemize}
\item
  $u_i\geq \delta_i$
  for all constraints of type $y_{i}-y_{i+1} \geq \delta_i$ in $Y$,
\item and $u_i<U_i$ and $u_i=\delta_i$ for all constraints
  $y_{i}-y_{i+1} = \delta_i$ in $Y$.
\end{itemize}
%
We say that $\alpha_f$ is a \emph{final abstract configuration
  respecting $Y$} whenever $q\in F$, $m_i=n_i$ for all $1\leq i\leq
d$, and $\alpha_f$ respects $Y$.


\subsection{Witnessing certificates}

While any concrete accepting run of an \laVASS gives rise to an
abstract configuration path ending in an accepting abstract
configuration, the converse does not hold. This motivates the
introduction of reachability and witnessing certificates, which are
special abstract configuration paths that carry further information
that eventually enables us to derive from a witnessing certificate a
concrete accepting run of an \laVASS.

A \emph{reachability certificate} is a tuple $(R,X,Y,L)$ such that
$R=\alpha_1\xrightarrow{t_1} \alpha_2 \xrightarrow{t_2} \cdots
\xrightarrow{t_{n-1}} \alpha_n$ is an abstract configuration path, and
$X,Y,L\colon \{1,\ldots,d\} \to \{1,\ldots,n\}$. Here, $X(i)$ and
$Y(i)$ indicate the position where the $x_i$-counter and $y_i$-counter
obtain a value different from $\bot$ for the first time. Moreover,
$L(i)$ is the position where a loop of length $\ell_i$ can be found.
Formally, $(R,X,Y,L)$ is required to have the following properties:
%
\begin{enumerate}[(a)]
\item $\alpha_1 = (q_0,\bot, \ldots, \bot,0,\ldots,0,
  \ell_1,\ldots,\ell_d)$ and if $\ell_i=\top$ then $\ell_j=\top$ for
  all $i<j\le d$;
\item $\lambda(x_i,t_{X(i)-1})=\wPP$ and $\lambda(y_i,t_{Y(i)-1})=\wTTPO$ for all $1\le i\le d$;
\item $\lambda(x_i,t_{j})=\wID$ for all $1\le j < X(i)-1$;
\item $\lambda(y_i,t_{j})=\wTT$ for all $1\le j < Y(i)-1$;
\item $X,Y,L$ are monotonic;
\item for all $1\le i\le d$, if $\ell_i\neq \top$ then
  \begin{itemize}
  \item $X(i)\le L(i)< Y(i)$; and
  \item there is a simple $\alpha_{L(i)}$-loop $\alpha_{L(i)}
    \xrightarrow{t_1'} \alpha'_2 \xrightarrow{t_2'} \cdots
    \xrightarrow{t_{\ell_i-1}} \alpha'_{\ell_i-1}
    \xrightarrow{t_{\ell_i}'} \alpha_{L(i)}$ of length $\ell_i$.
  \end{itemize}
\end{enumerate}
%
Those conditions can be interpreted as follows. Condition~(a) asserts
that the certificate starts in an initial abstract configuration. We
require that $\top$ monotonically propagates since since the absence of
a loop for counter $x_i$ implies that the remainder of a path is
short, hence we can afford to subsequently store actual counter
values and not residue classes. Conditions~(b), (c) and~(d) assert
that $X(i)$ and $Y(i)$ are the first position where the counters
$x_i,y_i$ hold value different from $\bot$. Condition~(e) states that
the counters $x_{i+1}$, $y_{i+1}$ do not carry a value different from
$\bot$ before the counters $x_{i}$ and $y_{i}$,
respectively. Condition~(f) implies that, if $\ell_i\neq \top$ then
between the first update for counter $x_i$ and the first update for
counter $y_i$ there is a position $L(i)$ where we can find a loop in
the abstract configurations of length $\ell_i$. Notice that if
$x_j=\bot$ or $y_j=\bot$ in $\alpha_{L(i)}$ then $x_j$ and $y_j$
remain to hold $\bot$ along this loop, i.e., this loop does not update
counters that have not been initialised already.

Given $R$, the set of $y$-constraints induced by $R$ is the smallest
set containing
%
\begin{itemize}
\item $y_{i}-y_{i+1}\geq \delta_i$, where $\delta_i \defeq
  X(i+1)-X(i)$ if there is a $j$ such that $X(i)\le L(j) < X(i+1)$;
  and
\item otherwise $y_{i}-y_{i+1} = \delta_i$, where $\delta_i \defeq
  X(i+1)-X(i)$,
\end{itemize}
for all $1\le i< d$ such that $\ell_i\neq\top$.

We introduce some further notation. Given a reachability certificate
$R$, we denote by $\pi(R)$ the run corresponding to $R$ in the
configuration graph of $V$, with the initial configuration
$(q_0,0,0,\ldots,0,0)$. Given indices $1\le i\le j\le n$, we denote by $R[i,j]$ the
segment $\alpha_i \xrightarrow{t_i} \alpha_{i+1} \cdots
\xrightarrow{t_{j-1}} \alpha_j$ of $R$, and by $R[i]\defeq \alpha_i$.
We say that $R$ is a \emph{witnessing certificate} if, for $a\le d$
being the largest index such that $\ell_a \neq \top$:
%
\begin{itemize}
\item $R[1,Y(a)]$ is a simple path and $n-Y(a)\le 2dM_{d+1}$;
\item $\alpha_n$ is a final abstract configuration respecting the set
  of induced $y$-constraints; and
\item $\val(\pi(R),x_a) \le \val(\pi(R),y_a)$.
\end{itemize}
%
Sometimes we will speak of witnessing certificates \emph{restricted}
to a set of counters. By that we mean a witnessing certificates where
the relevant Conditions~(a)--(f) are only required for that set of
counters.

Now we are ready to provide a proof for \Cref{prop:main}, that stated that language emptiness for restricted \laVASS can be decided in $\mathrm{NSPACE}( \abs V \cdot 2^{O(d)})$.
\begin{proof}[Proof of \Cref{prop:main}]
  Clearly, an abstract configuration can be stored in space $\abs V \cdot 2^{O(d)}$. An NEXPSPACE algorithm can hence
  non-deterministically choose an initial configuration and
  non-deterministically verify that it leads to a final abstract
  configuration along a path that is a witnessing certificate. To this
  end, the algorithm computes the set of induced $y$-constraints
  on-the-fly while guessing the reachability certificate, and verifies
  them in the last configuration. Note that the $y$-constraints can be stored in space $\abs V \cdot 2^{O(d)}$. Finally, the requirement
  $\val(\pi(R),x_a) \le \val(\pi(R),y_a)$ can also be verified in
  exponential space since we require that $R[1,Y(a)]$ is a simple path and $n-Y(a)\le 2M_{d+1}$.
\end{proof}
%
In the next section we argue the correctness of our algorithm by proving the following theorem:
\begin{theorem}\label{theorem:certificate iff non-empty language}
The language of a restricted \laVASS $V$ is non-empty if and only if there exists a witnessing certificate for $V$.
\end{theorem}
%



Given a restricted \laVASS $V$, we introduce the concept of \emph{abstract configurations}. 
\begin{definition}
An abstract configuration is a tuple 
$$\alpha=(q, m_1,n_1,m_2,n_2,\dots,m_{d},n_{d},u_1,u_2,\dots,u_{d-1},l_1,\dots,l_d),$$
where $q\in Q$, $m_1,n_1,\dots,m_d,n_d$, $u_1,\dots,u_{d-1}$ are natural numbers, $l_1,\dots,l_d\in\N\cup\{\top\}$, $m_i,n_i\leq (2\abs{Q})^{(4^{i-1})}$ and $u_i,l_i\leq\abs{Q}^{(4^{i-1})}$ or $l_i=\top$.
\end{definition}
The idea is that $m_i,n_i$ are storing the residue class modulo $l_i$ of counters $x_i,y_i$ respectively and we use the special symbol $\top$ to represent that we are storing actual values and not residue classes. The value of $u_i$ is storing the current difference between the value of counters $y_{i}$ and $y_{i+1}$. This difference is potentially unbounded so we only store its value if it is less than $l_i$ or record the fact that it is bigger or equal to $l_i$ (by having $u_i=\abs{Q}^{(4^{i-1})}$). Also, we do not allow the difference to become negative since $u_1,\dots,u_{d-1}$ are all natural numbers. 

The set of all abstract configurations is $A$. Note that this set if finite because $m_i,n_i,u_i,l_i$ are all bounded by $(2\abs{Q})^{4^{i-1}}$. We can define the following relation between abstract configurations, ${\xrightarrow{}}\subseteq A\times\mathcal{P}(\{x_1,y_1,\dots,x_{d},y_{d}\})\times A$ such that $\alpha\xrightarrow{I}\alpha'$ if and only if:
\begin{itemize}
    \item $\alpha=(q,m_1,n_1,\dots,u_{d-1},l_1,\dots,l_d)$ and $\alpha'=(q',m_1',n_1',\dots,u_{d-1}',l_1,\dots,l_d)$,
    \item $t=(q,a,q')\in\Delta$, for some $a\in\Sigma$,
    \item $\lambda(t)=(f_{x_1},f_{y_1},\dots,f_{x_d},f_{y_d})$,
    \item $f_{x_i} = id$ and $f_{y_i} = \timestwo$, for all $x_i,y_i\not\in I$,
    \item If $\l_i\neq\top$, $m_i' = f_{x_i}(m_i)\bmod l_i$ and $n_i'=f_{y_i}(n_i)\bmod l_i$ 
    \item If $\l_i=\top$, $m_i' = f_{x_i}(m_i)$ and $n_i'=f_{y_i}(n_i)$ and
    \item $$u_i' =
    \begin{cases*}
      min(2u_i+1,\abs{Q}^{(4^{i-1})}) & if $f_{y_{i}}=\timestwob,f_{y{i+1}}=\timestwo$ \\
      min(2u_i,\abs{Q}^{(4^{i-1})}) & if $f_{y_{i}}=f_{y_{i+1}}$ \\
      min(2u_i-1,\abs{Q}^{(4^{i-1})}) & if $f_{y_{i}}=\timestwo,f_{y{i+1}}=\timestwob,$ \\
    \end{cases*}
    $$
    for all $i\in\{1,\dots,d\}$
\end{itemize}


We write $\alpha\xrightarrow{I}^t \alpha'$ whenever $\alpha\xrightarrow{I}\alpha_1\xrightarrow{I}\dots \alpha_{t-1}\xrightarrow{I}\alpha'$ for some $\alpha_1,\dots,\alpha_t\in A$. Furthermore, we write $\alpha\xrightarrow{I}^* \alpha'$ whenever $\alpha\xrightarrow{I}^t \alpha'$ for some $t\geq 0$. Finally, we write $\alpha\xrightarrow{}^t\alpha'$ whenever $\alpha\xrightarrow{I}^t\alpha'$ for some $I\subseteq\{x_1,y_1,\dots,x_d,y_d\}$.  

Given two consecutive y-counters $y_i,y_{i+1}$ we define two different constraints: $y_{i}-y_{i+1}=t_i$ and $y_{i}-y_{i+1}\geq t_i$, where $t_i\geq 0$. We call these y-constraints. 

\begin{definition}
Let $Y$ be a set of y-constraints. We say that an abstract configuration $\alpha=(q,m_1,n_1,\dots,u_1,\dots,u_{d-1},l_1,\dots,l_d)$ respects $Y$ whenever $u_i=\abs{Q}^{4^{i-1}}$ or $u_i\geq t_i$ for all constraints of type $y_{i}-y_{i+1} \geq t_i$ in $Y$ and $u_i<\abs{Q}^{4^{i-1}}$ and $u_i=t_i$ for all constraints $y_{i}-y_{i+1} = t_i$ in $Y$. 
\end{definition}


We can now define a \emph{final abstract configuration}. 

\begin{definition}
Let $\alpha_f = (q,m_1,n_1\dots,u_1,\dots,u_{d-1},l_1,\dots,l_d)$ be an abstract configuration and $Y$ be a set of y-constraints. We say that $\alpha_f$ is a final abstract configuration with $Y$ if and only if $q\in F$, $m_i=n_i$ for all $1\leq i\leq d$ and $\alpha_f$ respects $Y$.
\end{definition}



We are now ready to define the certificate for the emptiness problem with $d$ pairs of counters. Let an $\emph{update pair}$ be a pair $(z,\alpha)$, where $z\in\{x_1,\dots,x_d,y_1,\dots,y_d\}$ is a counter ID and $\alpha$ is an abstract configuration. A $\emph{loop pair}$ is a pair $(l,\alpha)$, where $l\in\N\cup\{\top\}$ and $\alpha$ is an abstract configuration. The certificate $\zeta$ is a sequence of $2a$ update pairs $(z,c)$ interleaved with $a$ loop pairs $(l,c)$. If $a=d$, the certificate proves the existence of an accepting path on which all counters get updated at least once. Otherwise, if $a<d$ the certificate proves the existence of an accepting path on which only the counters that appear in an update pair get updated and the rest do not get updated (the value of these counters remains $0$).

Given such a certificate $\zeta$ for $V$, we now show how to verify the certificate. 
First, we do some checks on the order of the sequence of pairs in $\zeta$:
\begin{itemize}
    \item update pair $(x_i,\alpha)$ must come before update pair $(x_{i+1},\alpha')$,
    \item update pair $(y_i,\alpha)$ must come before update pair $(y_{i+1},\alpha')$,
    \item update pair $(x_i,\alpha)$ must come before update pair $(y_i,\alpha')$,
    \item between update pair $(x_i,\alpha)$ and update pair $(y_i,\alpha')$ there must be a loop pair $(l_i,\alpha'')$.
\end{itemize}

Let $\zeta[i]$ be the $i$-th update or loop pair in sequence $\zeta$ and $\zeta[i{:}j]$ be the sequence $\zeta[i],\zeta[i+1],\dots,\zeta[j]$, where $1\leq i\leq j\leq 3a$. Also, for a sequence of update and loop pairs $\zeta'$ let $I[\zeta']$ be the set of all counters that are mentioned in an update pair in $\zeta'$. 

We now check for all $i \in\{1,\dots,3a\}$ that:

\begin{enumerate}
    \item \label{check1} $\alpha_0\xrightarrow{\emptyset}^*\alpha_1$, where $\alpha_0$ is the abstract configuration representing the initial configuration of $V$, $\alpha_0=(q_0,0,0,\dots,l_1,\dots,l_a)$ and $\zeta[1]=(x_1,\alpha_1)$
    \item \label{check2} $\alpha\xrightarrow{I[C[1{:}i]]}^* \alpha'$, where $\zeta[i] = (\cdot, \alpha)$ and $\zeta[i+1] = (\cdot,\alpha')$ (if $i<3a$). Also, if $\zeta[i]$ is an update pair $(x_i,\alpha)$ or $(y_i,\alpha)$, we check that the first transition of $\alpha\xrightarrow{I[\zeta[1{:}i]]}^* \alpha'$ is labelled with $++$ for $x_i$ or $\timestwob$ for $y_i$ respectively. 
    \item \label{check3} if $\zeta[i]$ is a loop pair $(l,\alpha)$ and $l\neq\top$, then $\alpha\xrightarrow{I[\zeta[1{:}i]]}^{l} \alpha$.
\end{enumerate}

These paths are easy to verify because the set of all abstract configurations $A$ is finite. 


We are also building $Y$, a set of y-constraints. For all $i\in\{1,\dots,a-1\}$, by the previously identified paths, we know that $\alpha\xrightarrow{}^{t_i} \alpha'$, where $(x_i,\alpha),(x_{i+1},\alpha')\in \zeta$, for some $t_i\geq 0$. If there is no loop pair in $\zeta$ between pairs $(x_i,\alpha)$ and $(x_{i+1},\alpha')$ we add constraint $y_{i}-y_{i+1} = t_i$ to $Y$. Otherwise, we add constraint $y_{i}-y_{i+1}\geq t_i$ to $Y$.  

Finally, let $(y_a,\alpha)$ be $\zeta[3a]$. We know that there exists update pair $(x_a,\alpha')\in\zeta$ and that $\alpha'\xrightarrow{\{x_1,y_1,\dots,x_a,y_{a-1}\}}^t\alpha$, for some $t\geq 0$. We must check that $\alpha\xrightarrow{\{x_1,\dots,y_a\}}^{t'} \alpha_f$, where $\alpha_f$ is a final abstract configuration with $Y$ and the final value of $n_a$ is greater or equal to $t$. 


We are now ready to introduce the central theorem of this paper:
\begin{theorem}
Let $V$ be a restricted \laVASS. $L(V)\neq\emptyset$ if and only if there exists a certificate $\zeta$ that can be verified.
\end{theorem}
\begin{proof}

$(=>):$

If $L(V)\neq\emptyset$ it means that there must exist an accepting path $\pi$. So, there exists a sequence of configurations $P = \{c_i\}_{0\leq i\leq r}$ and alphabet symbols $\{a_i\in\Sigma\}_{0<i\leq r}$ s.t. $c_i\xrightarrow{a_i}c_{i+1}$ for all $0\leq i< r$, $c_0=(q_0,0,\dots,0)$ and $c_r$ is a final configuration. 

Let $0\leq a\leq d$ be the number of x-counters that get updated at least once (their final value is bigger than 0). We can define $P[x_i]$ to be the configuration before we see the first update for counter $x_i$ and $P[y_i]$ to be  the configuration before we see the first update for counter $y_i$ in $P$. Finally, let $c_f=c_r$. Let $R$ be the sequence of configurations $c_f$ and $P[x_i],P[y_i]$ for all $i\in\{1,\dots,a\}$ that identify.

Our aim is to turn the sequence $R$ into a certificate $\zeta$. For this, we need a way to turn a configuration $c$ into an abstract configuration $\alpha$. Given a natural number $a$ and a restricted \laVASS $V$, let $f_a^V:C\times (\N\cup\{\top\})^a\xrightarrow{}A$ be a function that takes a configuration and loop variables $l_1,\dots,l_a$ and returns an abstract configuration. 

\begin{gather*}
f_a^V((q,m_1,n_1,\dots,m_d,n_d),l_1,\dots,l_a)=\\
(q,m_1\circ l_1,n_1\circ l_1,\dots,m_a\circ l_a,n_1\circ l_a,m_{a+1},n_{a+1},\dots,m_d,n_d,\\
min(n_1-n_2,\abs{Q}+1),\dots,min(n_{d-1}-n_{d},\abs{Q}^{4^{d-1}}+1),l_1,\dots,l_a,\top,\dots,\top),\\
\text{where } m\circ l=
\begin{cases*}
m\bmod l & if $l\in\N$\\
m & if $l=\top$
\end{cases*}
\end{gather*}

We say that $f_a^V((q,m_1,n_1,\dots,m_d,n_d),l_1,\dots,l_a)$ is well-defined if 
\begin{itemize}
    \item $l_i=\top$ or $l_i\leq\abs{Q}^{4^{i-1}}$,
    \item $m_i,n_i\leq (2\abs{Q})^{4^{i-1}}$ if $l_i=\top$, for all $i\in\{1,\dots,a\}$ and
    \item $n_{i}-n_{i+1}\geq 0$, for all $i\in\{1,\dots,d-1\}$.
\end{itemize}

The function $f_a^V$ assigns the value $\top$ to $l_{a+1},\dots,l_d$ because the accepting path $\pi$ does not update the values of counters $x_{a+1},y_{a+1},\dots,x_d,y_d$ and $m_{a+1}=n_{a+1}=\dots,m_d=n_d=0$ in all configurations of $\pi$, so the values of $l_{a+1},\dots,l_d$ do not matter. 

The next step is to lift a path in $V$ on configurations to a path on abstract configurations. 

\begin{lemma}
\label{lm:path}
Let $c,c'\in C$ be configurations, $0\leq a\leq d$ and $l_1,\dots,l_a\in\N\cup\{\top\}$. If $c\xrightarrow{} c'$, $f_a^V(c,l_1,\dots,l_a)$ and $f_a^V(c',l_1,\dots,l_a)$ are well-defined then $f_a(c,l_1,\dots,l_a)\xrightarrow{I} f_a^V(c',l_1,\dots,l_a)$, where $I$ is the set of counter IDs that have value zero in $c'$.
\end{lemma}

\begin{proof}
Follows from the definition of $\xrightarrow{}$ and $f_a^V$.
\end{proof}

\begin{theorem}
\label{th:ls}
Let $V$ be a restricted \laVASS, $0\leq a\leq d$ and $P$ be a sequence of configurations of $V$ s.t. the first configuration is the initial one, the last configuration is a final configuration $(q,m_1,m_1,\dots,m_a,m_a,0,\dots,0)$, with $m_i\geq m_{i+1},\forall i\in\{1,\dots,a-1\}$ and $\forall c,c'$ consecutive configurations in $P$, $c\xrightarrow{}c'$ in $V$. There exist $l_1\dots,l_a$ s.t.
\begin{itemize}
    \item $\forall c\in P,f_a^V(c,l_1,\dots,l_a)$ is well-defined,
    \item $l_i=\top$ or between $P[x_i]$ and $P[y_i]$ there exists a configuration $P[L_i]\in P$ s.t. $\alpha\xrightarrow{I}^{l_i}\alpha$, where $\alpha=f_a(P[L_i],l_1,\dots,l_a)$ and $I$ is the set of counter IDs that do not have value zero in $P[y_i]$, for all $i\in\{1,\dots,a\}$.
\end{itemize}
\end{theorem}

\begin{proof}

Since we know that the final configuration of $P$ is $(q,m_1,m_1,\dots,m_a,m_a,0,\dots,0)$, with $m_i\geq m_{i+1},\forall i\in\{1,\dots,a-1\}$, using the definition of a restricted \laVASS we can conclude that for all $(q,m_1',n_1',\dots,m_a',n_a',0,\dots,0)\in P$, $m_i'\geq m_{i+1}'$ and $n_i'\geq n_{i+1}'$.

Now, consider the path $P[x_1]\xrightarrow{}^* P[y_1]$. Assume that $P[x_1]\xrightarrow{}^t P[y_1]$ and $t\leq\abs{Q}$. By the definition of a restricted \laVASS we know that after $k$ transitions from $P[y_1]$:
\begin{gather*}
    \text{The value of counter } x_1<\abs{Q}+k\\
    \text{The value of counter } y_1\geq 2^k
\end{gather*}

So, for $k=\abs{Q}$, we get that the value of counter $y_1$ is greater than the value of counter $x_1$ and counter $x_1$ can not recover this difference. This implies that the maximum value of counters $x_1$ and $y_1$ is less than $2\abs{Q}$, thus the last configuration of $P$ must be $c_f=(q_f,m_1',m_1',\dots,m_a',m_a',0,\dots,0)$ with $q_f\in F$ and $m_1'<2\abs{Q}$. So, if we pick $l_1=\dots=l_a=\top$ we have that $\forall c\in P$, $f_a(c,l_1,\dots,l_a)$ is well-defined because any configuration $c\in P$ has only counter values that are less than $2\abs{Q}$ and the second condition of the theorem is trivially satisfied. 

If $P[x_1]\xrightarrow{}^t P[y_1]$ and $t>\abs{Q}$ then the path must contain a loop, so there exist $l_1\leq\abs{Q}$ and configuration $P[L_1]=(q,m_1,0,\dots)$ between $P[x_1]$ and $P[y_1]$ s.t. $P[L_1]\xrightarrow{}^{l_1}(q,m_1+l_1,0,\dots)$. We can state that $\forall c\in P$, $f_a^V(c,l_1,1,\dots,1)$ is well-defined because picking 1 for the values of $l_2,\dots,l_a$ clearly satisfies the condition that $1\leq\abs{Q}^{4^{i-1}}$. Also, there exists $P[L_1]\in P$ s.t. $\alpha\xrightarrow{I}^{l_1}\alpha$, where $\alpha=f_a^V(P[L_1],l_1,1,\dots,1)$ and $I$ is the set of counter IDs that do not have value zero in $P[y_1]$.

Assume now that there exist $l_1,\dots,l_{k-1}$ s.t. for all $i\in\{1,\dots,k-1\}$:
\begin{itemize}
    \item $l_i\neq\top,$
    \item $\forall c\in P$, $f_a^V(c,l_1,\dots,l_{k-1},1,\dots,1)$ is well-defined and
    \item between $P[x_i]$ and $P[y_i]$ there exists a configuration $P[L_i]\in P$ s.t. $\alpha\xrightarrow{I}^{l_i}\alpha$, where $\alpha=f_a^V(P[L_i],l_1,\dots,l_{k-1},1,\dots,1)$ and $I$ is the set of counter IDs that do not have value zero in $P[y_i]$.
\end{itemize}

We aim to find a suitable value for $l_k$. Consider the case in which $P[L_{k-1}]$ comes after $P[x_k]$ in $P$. In this case, one can simply pick $l_k=l_{k-1}$ and trivially obtain that there exists $P[L_k]=P[L_{k-1}]$ s.t. $\alpha\xrightarrow{I}\alpha$, where $\alpha = f_a^V(P[L_k],l_1,\dots,l_k,1\dots,1)$ and $I$ is the set of IDs that do not have value zero in $P[x_k]$. The previous loop on abstract states is preserved because the value of $n_k=0$ for all configurations before $P[y_k]$ by definition and $m_k\bmod l_k = (m_k+l_k)\bmod l_k$. Now, we need to consider cases in which $P[L_{k-1}]$ comes before $P[x_k]$ in $P$.

We distinguish two more cases.

Case 1. $P[x_k]\xrightarrow[]{}^t P[y_k]$ with $t\leq\abs{Q}^{4^{k-1}}$. Again, we know that after $\abs{Q}^{4^{k-1}}$ transitions from $P[y_k]$ the value of counter $y_k$ will exceed the value of counter $x_k$ and counter $x_k$ will not be able to recover this difference. ($2\abs{Q}^{4^{k-1}}<2^{\abs{Q}^{4^{k-1}}}$, for all $k$). So, the maximum value of counter $y_k$ is less then $2\abs{Q}^{4^{k-1}}$ which implies that for all $c\in P$, $m_k,n_k,\dots,m_a,n_a\leq 2\abs{Q}^{4^{k-1}}$. Thus, if we pick $l_k=l_{k+1}=\dots,l_a=\top$ we have that $\forall c\in P$, $f_a^V(c,l_1,\dots,l_a)$ is well-defined. Moreover we know from our assumption that between $P[x_i]$ and $P[y_i]$ there exists a configuration $P[L_i]\in P$ s.t. $\alpha\xrightarrow{I}^{l_i}\alpha$, where $\alpha=f_a^V(P[L_i],l_1,\dots,l_{k-1},1,\dots,1)$ and $I$ is the set of counter IDs that do not have value zero in $P[y_i]$, for all $i\in\{1,\dots,k-1\}$ and that $P[k-j]$ comes before $P[x_k]$ for all $j>0$. So, all configurations before $P[x_k]$ have value zero for $m_k,n_k,\dots,m_a,n_a$ and thus we can safely state that between $P[x_i]$ and $P[y_i]$ there exists a configuration $P[L_i]\in P$ s.t. $\alpha\xrightarrow{I}^{l_i}\alpha$, where $\alpha=f_a^V(P[L_i],l_1,\dots,l_{a})$ and $I$ is the set of counter IDs that do not have value zero in $P[y_i]$, for all $i\in\{1,\dots,k-1\}$.

Case 2. $P[x_k]\xrightarrow{}^t P[y_k]$ with $t>\abs{Q}^{4^{k-1}}$. Let us count the number of different outputs that we can get for $f_a^V(c,l_1,\dots,l_{k-1},1\dots,1)$, where $c$ can be any configuration that may appear on the path $P[x_k]\xrightarrow{}^t P[y_k]$. We know that $c$ is of the form $(q,m_1,n_1,\dots,m_a,n_a,0,\dots,0)$ and $n_k=0$ by the definition of $P[y_k]$. However, we also know that if $c\in P$, $n_i\geq n_{i+1}$, $\forall i\in\{1,\dots,d-1\}$, so $n_k=\dots=n_a=0$. This implies that for any $\alpha=f_a^V(c,l_1,\dots,l_{k-1},1\dots,1)$, $u_k=u_{k+1}=\dots=u_{a-1}=0$. Also, $m_k=n_k=\dots=m_{a}=n_a=0$ because $l_k=\dots=l_a=1$. Thus, the number of different outputs of $f_a^V(c,l_1,\dots,l_{k-1},1,\dots,1)$ is:
\begin{align*}
&\abs{Q}\times(\prod_{1\leq i<k}l_i)^2\times\prod_{1\leq i<k}(\abs{Q}^{4^{i-1}})\\
\leq &\abs{Q}\times(\prod_{1\leq i<k} \abs{Q}^{4^{i-1}})^3\\
=&\abs{Q}\times(\abs{Q}^{\sum_{0\leq i<k-1}4^i})^3\\
=&\abs{Q}^{4^{k-1}}
\end{align*}

So, by the pigeonhole principle, we know that there exists natural number $l_k\leq\abs{Q}^{4^{k-1}}$ and configuration $P[L_k]$ between $P[x_k]$ and $P[y_k]$ s.t. $\alpha\xrightarrow{I}^{l_k}\alpha$, where $\alpha=f_a^V(P[L_k],l_1,\dots,l_{k-1},1,\dots,1)$ and $I$ is the set of counter IDs that do not have value zero in $P[y_k]$. However, since we know that in $P[L_k]$ the value of $n_k$ is zero and taking $l_k$ transitions changes the value of $m_k$ by $l_k$ we can also state that $\alpha'\xrightarrow{I}^{l_k}\alpha'$, where $\alpha'=f_a^V(P[L_k],l_1,\dots,l_k,1,\dots,1)$.

\end{proof}

So, we started with an accepting path $\pi$ and we saw how to extract the relevant configurations from $\pi$ ($c_f$ and $P[x_i],P[y_i]$). From here, we can relabel the counters such that either $P[x_i]=P[x_{i+1}]$ or $P[x_i]$ appears before $P[x_{i+1}]$ and either $P[y_{i}]=P[y_{i+1}]$ or $P[y_i]$ appears before $P[y_{i+1}]$ in $P$ for all $i$. Since $P$ describes an accepting path, we know that $P[x_i]$ is either equal to $P[y_{i}]$ or $P[x_{i}]$ appears before $P[y_{i}]$. Otherwise, if counter $y_i$ is updated before counter $x_i$, counter $x_i$ will never be able to recover the difference and thus the value of $x_i$ cannot be equal to the value of $y_i$. Now, we can apply Theorem \ref{th:ls} to obtain $l_1,\dots,l_a$, so that we can apply $f_a^V(\cdot,l_1,\dots,l_a)$ to all configurations of $P$ to obtain a certificate $\zeta$. Update pairs $(x_i,\alpha)$ and $(y_i,\alpha)$ are  $(x_i,f^V_a(P[x_i],l_1,\dots,l_a))$ and $(y_i,f^V_a(P[y_i],l_1,\dots,y_a))$ respectively and loop pair $(l_i,\alpha)$ is $(l_i,f^V_a(P[L_i],l_1,\dots,l_a))$ if $l_i\neq\top$ or $(x_i, f^V_a(P[x_i],l_1,\dots,l_a))$ if $l_i=\top$. Bellow is a high-level description of the process for obtaining a certificate $\zeta$ from an accepting path $\pi$.

$$
\pi\xrightarrow[\text{important cfs.}]{\text{extract}}P\xrightarrow[\text{counters}]{\text{relabel}}P\xrightarrow{\text{Theorem \ref{th:ls}}}P,l_1,\dots,l_a,P[L_1],\dots,P[L_k]\xrightarrow[f^V_a]{\text{apply}}\zeta
$$

We now prove that $\zeta$ can be verified. The order of $\zeta$ is verifiable because of the relabelling and the way we choose $P[L_1],\dots,P[L_a]$ (either $P[L_i]$ is between $P[x_i]$ and $P[y_i]$ or $P[L_i]=P[x_i]$, for all $i\in\{1,\dots,a\}$). Next, we know that checks \ref{check1} and \ref{check2} are satisfied by of Lemma \ref{lm:path} and by our definition of $P[x_i],P[y_i]$ and check $\ref{check3}$ is satisfied by Theorem $\ref{th:ls}$. 

Now, assume that $c_f=(q_f,m_1,n_1,\dots,m_a,n_a,0,\dots,0)$. Since $c_f$ is a final configuration we know that $m_i=n_i$ for all $i\in\{1,\dots,a\}$. Consider the set of y-constraints that we are building during the verification phase. For all $i\in\{1,\dots,a-1\}$ we have $\alpha\xrightarrow{}^{t_i}\alpha'$, where $(x_i,\alpha),(x_{i+1},\alpha')\in\zeta$, for some $t_i\geq 0$. If $l_i=\top$ we know that $P[x_i]\xrightarrow{}^t P[y_i]$ with $t\leq\abs{Q}^{4^{i-1}}$, so $m_i-m_{i+1}<\abs{Q}^{4^{i-1}}$, so $u_i=n_i-n_{i+1}<\abs{Q}^{4^{i-1}}$, and $u_i=t_i$. If $l_i\neq\top$ we know that $P[x_i]\xrightarrow{}^t P[y_i]$ with $t>\abs{Q}^{4^{i-1}}$, so $m_i-m_{i+1}>\abs{Q}^{4^{i-1}}$, so $u_i=\abs{Q}^{4^{i-1}}$. In both cases we have that $f_a^V(c_f,l_1,\dots,l_a)$ is a final configurations with $Y$. Finally, since $m_a=n_a$ we must have that $n_a$ is equal to $m_a$ which is equal to the length of the path from $P[x_a]$ to $c_f$. 

$(<=):$

Assume a verifiable certificate $\zeta=\langle(x_1,\alpha_1),\dots,(y_a,\alpha_{3a})\rangle$ exists. Since $\zeta$ is verifiable we know that $\alpha_0\xrightarrow{\emptyset}^*\alpha_1\xrightarrow{\{x_1\}}^*\dots\xrightarrow{}^*\alpha_{3a}\xrightarrow{\{x_1,\dots,y_a\}}^*\alpha_f$, where $\alpha_0$ is the abstract configuration $(q_0,0,\dots,0)$ and $\alpha_f$ is a final abstract configuration. 

This implies that there exists a path $\pi$ in $V$ from the initial configuration $(q_0,0,\dots,0)$ to a configuration $(q_f,m_1,n_1,\dots,m_a,n_a,0,\dots,0)$, where $q_f\in F$ in which:
\begin{itemize}
    \item $m_i\bmod l_i = n_i\bmod l_i$ (or $m_i=n_i$ if $l_i=\top$) (because of the definition of $\xrightarrow{}$),
    \item $n_a\geq m_a$ (because of the final certificate check) and
    \item $n_{i}-n_{i+1}\geq u_i$ (because of the definition of $\xrightarrow{}$), for all $i\in\{1,\dots,a\}$.
\end{itemize}

Moreover, by the way we construct $Y$ and the fact that $\alpha_f$ is a final abstract configuration with respect to $Y$, we know that $c\xrightarrow{I}^{t_i}c'$, where $(x_i,c),(x_{i+1},c')\in \zeta$, for some $I\subseteq\{x_1,\dots,y_a\}$, $u_i\geq t_i$ (since $\alpha_f$ respects $Y$), so $m_{i}-m_{i+1}\leq u_i$, for all $1\leq i\leq a$.


Thus, we can conclude that $m_i\leq n_i$ for all $1\leq i\leq a$. Putting it all together we know that there exists a path $\pi$ in $V$ that goes from the initial configuration to a configuration $(q_f,m_1,n_1,\dots,m_a,n_a,0,\dots,0)$, where $q_f\in F$ and 

\begin{align*}
    &n_1 = m_1 + l_1\times k_1\text{ or } n_1 = m_1\text{ if } l_1=\top\\
    &n_2 = m_2 + l_2\times k_2\text{ or } n_2 = m_2\text{ if } l_2=\top\\
    &\dots\\
    &n_a = m_a + l_d\times k_a\text{ or } n_a = m_a\text{ if } l_a=\top,
\end{align*}

for some $k_1,\dots,k_a$ natural numbers.

If $l_1=\dots=l_a=\top$ we are done because if $\zeta$ is verifiable it means that we can obtain a path $\pi$ in $V$ from $\zeta$ that starts with the initial configuration and ends in a final configuration. $(m_1=n_1,\dots,m_a=n_a)$. 

Otherwise, assume $l_{i+1}=l_{i+2}=\dots=l_a=\top$, for some $1<i\leq a$. So, $\pi$ ends in a final configuration $(q_f,m_1,n_1,\dots,m_a,n_a)$ in which:

\begin{align*}
    &n_1 = m_1 + l_1\times k_1\\
    &n_2 = m_2 + l_2\times k_2\\
    &\dots\\
    &n_i = m_i + l_i\times k_i\\
    &n_{i+1} = m_{i+1}\\
    &\dots\\
    &n_a = m_a   
\end{align*}

We are interested in adding some extra spins to $\pi$ in order to increase the values of counters $m_1,\dots,m_i$ to get full equality. Concretely, if $m_{i}\neq n_{i}$ we prove that there exists a path $\pi'$ that uses the loop given by the loop pair $(l_{i},\alpha_l)$ to get to a configuration in which $m_{i}=n_{i}$ while preserving the congruence relations on $m_1,n_1,\dots,m_{i-1},n_{i-1}$. 

We know that between update pairs $(x_i,\alpha)$ and $(y_i,\alpha')$ there exists a loop pair $(l_i,\alpha_l)$. Between the loop pair $(l_i,\alpha_l)$ and the update pair $(y_i,\alpha')$ we might have many update pairs for some counters, but we do not have to consider them since adding spins around loop $L_i$ does not change the value of those counters. However, we do have to consider the update pairs between $(x_i,\alpha)$ and $(l_i,\alpha_l)$. 

Case 1, update pair $(x_{i+1},\alpha'')$ is between $(x_i,\alpha)$ and $(l_i,\alpha_l)$. Note that this case implies that $i<a$. Then, we know that $\alpha\xrightarrow{}^{t_i} \alpha''$, so $m_i-m_{i+1} = t_i$. Moreover, $n_{i}-n_{i+1}=t_i$ because $\pi$ goes to a final abstract configuration that respects $Y$ ($u_i=t_i$). Finally, because $n_{i+1}=m_{i+1}$ we can conclude that $n_{i} = m_i$, thus $\pi'=\pi$. 

Case 2, update pair $(x_{i+1},\alpha'')$ is not between $(x_i,\alpha)$ and $(l_i,\alpha_l)$. We know that $\alpha_l\xrightarrow{}^{l_i}\alpha_l$, so let $L$ be the loop defined by $\alpha_l\xrightarrow{}^{l_i}\alpha_l$.

Let $\pi'$ be the path we get after we add $k_i$ spins around around $L$. In other words, $\pi'$ is the path that we obtain from
$$\alpha_0\xrightarrow{}^*\dots\xrightarrow{}^*(\alpha_l\xrightarrow{}^{l_i}\alpha_l)^{k_i}\xrightarrow{}^*\dots\xrightarrow{}^*\alpha_f.$$
After $\pi'$ we reach a configuration $(q_f,m_1',n_1',\dots,m_a',n_a',0,\dots,0)$, where $n_i'=n_{i},m_{i+1}'=m_{i+1},n_{i+1}'=n_{i+1},\dots,m_{a}'=m_{a'},n_{a}'=n_a$ because the extra spins around $L$ cannot change the values of the counters that get updated after $L$. Also, $m_{i}' = m_{i}+k_il_i$, so $m_i'=n_i'$. Lastly, because the path ends with the same final abstract configuration, we know that:

\begin{align*}
    &n_1' = m_1' + l_1\times k_1'\\
    &n_2' = m_2' + l_2\times k_2'\\
    &\dots\\
    &n_{i-1}' = m_{i-1}' + l_{i-1}\times k_{i-1}'\\
    &n_i' = m_i'\\
    &n_{i+1}' = m_{i+1}'\\
    &\dots\\
    &n_a' = m_a'   
\end{align*}

So, we can conclude that there exists a path $\pi_f$ that uses loop pairs $(l_1,\alpha_1),\dots(l_{i-1},\alpha_{i-1})$ that starts with the initial configuration and ends with an accepting configuration $(q_f,m_1,n_1,\dots,m_a,n_a,0,\dots,0)$, where $m_1=n_1,\dots,m_a=n_a$. 

\end{proof}
